\section{Additional empirical results}\label{app:emp}
% =============================================================================================
% Tables: tab:e1-size-real (generated e1_size_real.tex), tab:e1-nonstat (manual_e1_nonstat_app.tex,
% narrower copy of the generated e1_nonstat.tex, same numbers), tab:e3-perticker-null (generated),
% tab:e3-likelihood-ag (generated), tab:e3-witness-app (manual_e3_witness_app.tex, derived from
% results/e3/e3_witness.csv). Scripts: paper/scripts/*.py. Derived quantities quoted in
% Section 7 are documented in app:emp-derived.
% =============================================================================================

This appendix collects the tables behind \cref{sec:emp} that are not in the main text and
documents the quantities that we derived from the stored results rather than read from a
generated table.

\subsection{Real-data placebo}\label{app:emp-real}

\Cref{tab:e1-size-real} reports E1(b): the 100 panel tickers are split at random into two halves
of 50, and the full pipeline tests one half against the other on the same dates, 50 times within
\TEST{} and 50 times within \TRAIN{}. In the \TRAIN{} splits, the relative tests are also run with
$\Rnew$ the real \TEST{} windows, half B in the role of the generator and half A in the role of
$\Rpast$, a generator equal in law to the past up to composition. The limits of this placebo are
discussed in \cref{sec:emp-placebo}: the Driscoll--Kraay rows are conservative by construction
because the halves share their dates, the halves need not have the same pooled window law, and
the splits are dependent and conditional on one market path, so we show rates without intervals.

Within \TEST{} the kernel tests with the wild bootstrap reject in 0.340 (MMD) and 0.400
(\SigMMD{}) of the splits at 5\%, and the time-block permutation p-value of the \WC{} test in
0.200. Within \TRAIN{} the wild-bootstrap rates are 0.020 for MMD and 0.200 for \SigMMD{}. The
anchored relative test does not reject in any \TRAIN{} split with its Driscoll--Kraay p-value and
in 0.020 of them with its permutation p-value. The two-critic test with the paired bootstrap
rejects in 0.240, which illustrates once more that it is anti-conservative at the boundary
(\cref{prop:boundary-size}).

\begin{table}[tbp]
  \centering
  \caption{Real-data placebo (E1(b)): rejection rates at $\alpha\in\{0.01,0.05,0.10\}$ when one
    random half of the 100 panel tickers is tested against the other half on the same dates; 50
    random splits within \TEST{} and 50 within \TRAIN{}; $L=20$. Relative tests (\TRAIN{} splits):
    $\Rnew$ = all real \TEST{} windows, $G$ = half B, $\Rpast$ = half A. Learned and relative
    tests use the time-block split; $^{\star}$marks the protocol's p-value (Driscoll--Kraay for
    learned tests, wild bootstrap for kernel tests). Fixed-$b$ rows: the Driscoll--Kraay
    statistic against fixed-$b$ references \citep{kiefer_2005_new}, conservative single-series
    reference (cons.) or two-sample equal-variance reference (eq.\ var.). Rates only: the splits
    are dependent and conditional on one realised market path, so binomial intervals do not
    apply.}
  \label{tab:e1-size-real}
  \small
  \fittable{e1_size_real}
\end{table}

\subsection{Past against future}\label{app:emp-nonstat}

\Cref{tab:e1-nonstat} reports E1(c), the absolute tests of the 3000 real \TEST{} windows against
the 6200 real \TRAIN{} windows with the E3 pipeline, over ten test seeds (critic initialisation,
block assignment, kernel subsample). Next to each p-value type it gives the rejection rate of the
same comparison in the stationary simulated placebo (E1(a), cell ``joint'', $\Rnew^*$ against
$\Rpast^*$, 200 replications), which is the size of that p-value for this comparison when the
two periods have the same law. The protocol p-values (Driscoll--Kraay, wild bootstrap) reject
with rates between 0.055 and 0.130 in that placebo; the permutation p-values and KS reject in 0.355 to 1.000 of its replications. On
the real data, the valid p-values (protocol and fixed-$b$) give $\pmed$ between 0.071 (\LRC{}) and
1.000 (\SigMMD{}), so the window tests do not reject \TRAIN{} against \TEST{} windows of
normalised returns at 5\%. With about 15 \TEST{} evaluation dates they have little power for this
comparison, and because the Driscoll--Kraay variance is not the same when both samples are factor
panels as when the synthetic tickers are independent, these p-values are not comparable with the
absolute p-values of E3 (\cref{sec:emp-placebo}). The \WC{}
statistic varies across the ten seeds from $-0.273$ to $0.440$, and is negative in two of them
(\cref{app:emp-derived}); which \TRAIN{} blocks train the critic matters
(\cref{rem:nonstat-past}).

\begin{table}[tbp]
  \centering
  \caption{Real \TEST{} windows against real \TRAIN{} windows (E1(c)): 3000 against 6200 windows,
    $L=20$, dependent protocol, ten test seeds. Median statistic (\WC{}: critic mean difference;
    \LRC{}: mean logit difference; C2ST: evaluation accuracy; MMD and \SigMMD{}: unbiased
    $\MMD^2$; KS: $D$), median p-value, $\pmed$ and share of seeds with $p\le0.05$. Last column:
    rejection rate at 5\% of the same comparison in the stationary simulated placebo (E1(a),
    joint cell, $\Rnew^*$ against $\Rpast^*$, 200 replications, Wilson 95\% interval).
    $^{\star}$Protocol p-value. Same numbers as the generated table \texttt{e1\_nonstat.tex}
    without its units and median $\zhac$ columns.}
  \label{tab:e1-nonstat}
  \small
  \fittable{manual_e1_nonstat_app}
\end{table}

\subsection{Per-ticker null and likelihood comparisons}\label{app:emp-perticker}

\Cref{tab:e3-perticker-null} gives the size of the nominal per-ticker p-values ($L=5$) under the
per-ticker \GARCHt{} null used to calibrate them: for each ticker, ten pairs of independent
\TEST{}-length paths from the fitted end-of-\TRAIN{} state, tested exactly like the real data. The
nominal p-values are oversized at every level, more for MMD than for the \WC{} test, and the last
column gives the nominal p-value that corresponds to 5\% after calibration. The draws of one
ticker share its law, so the Wilson intervals are approximate.

\Cref{tab:e3-likelihood-ag} gives the pairwise Amisano--Giacomini comparisons
\citep{amisano2007comparing}: per ticker, the HAC-studentised mean of the daily log-score
differences $d_t$ over the 615 \TEST{} days, without parameter penalty; pooled,
the cross-sectional mean of $d_t$ with a Driscoll--Kraay standard error. \GARCHt{} improves on GBM
by 0.1748 and on Merton by 0.0326 nats per day and per ticker on average, with pooled $z$ of 15.72
and 7.98; Merton improves on GBM by 0.1422. No ticker favours the simpler model significantly in
any pair after the Benjamini--Hochberg correction.

\begin{table}[tbp]
  \centering
  \caption{Size of the nominal per-ticker p-values under the per-ticker \GARCHt{} null (E3,
    $L=5$, 123 windows per sample): 10 independent pairs of \TEST{}-length paths per ticker, 1000
    in all; Wilson 95\% intervals (approximate: pairs of one ticker share its law). Last column:
    5\% quantile of the null p-values, the nominal p-value that corresponds to 5\% after
    calibration.}
  \label{tab:e3-perticker-null}
  \small
  \fittable{e3_perticker_null}
\end{table}

\begin{table}[tbp]
  \centering
  \caption{Amisano--Giacomini comparisons of the predictive log-scores on \TEST{} (E3, 615 days,
    100 tickers). $d_t$: daily log-score of model a minus model b. Pooled: cross-sectional mean of
    $d_t$, Driscoll--Kraay $z$ and p-value. Per ticker: HAC-studentised mean of $d_t$; ``a
    better'' and ``b better'': tickers with a significant difference in that direction after
    Benjamini--Hochberg at 5\%; last column: tickers with $z>0$.}
  \label{tab:e3-likelihood-ag}
  \small
  \fittable{e3_likelihood_ag}
\end{table}

\subsection{Witness diagnostics and other derived quantities}\label{app:emp-derived}

Some statements in \cref{sec:emp} rest on computations on the stored per-replication results
rather than on a generated table. We list them here with their sources (all under
\texttt{results/}), so that each can be checked.

\emph{Witness} (\cref{sec:emp-witness}). \Cref{tab:e3-witness-app} is computed from the critic
values stored by the witness part of E3 (\texttt{e3/e3\_witness.csv}): for each generator, the
seed-0 critic evaluated on its 1500 real and 1500 synthetic evaluation windows. The ``share of
$T$ from realised volatility'' replaces the critic by its mean within 20 equal-count bins of the
window log realised volatility (pooled over both samples) and divides the resulting
real-minus-synthetic difference by $T$. Realised volatility is the root mean square of the
normalised returns of the window, so 1 is the \TRAIN{} standard deviation of the ticker.

\emph{Standard errors} (\cref{sec:emp-power}). The ratio 7.7 is $0.245/0.032$ from
\cref{tab:e3-power}. The ratios of the real-data standard error of $\Danc$ (median over the five
seeds) to the D1 and D2 values are stored per generator in \texttt{e3/e3\_verdicts.csv}: 6.07, 5.16,
4.81, 4.91 and 4.53 times the D1 value and 0.79, 0.67, 0.62, 0.64 and 0.59 times the D2 value for
GBM, Merton, \GARCHt{}, Heston and the bootstrap.

\emph{Orderings across seeds} (\cref{sec:emp-tests}). From the per-seed statistics in
\texttt{e3/e3\_pooled\_raw.csv}: in each of the five seeds the absolute \WC{} statistic orders
the generators GBM, Merton, Heston, bootstrap, \GARCHt{}, and the anchored statistic orders the
four parametric generators GBM, Merton, Heston, \GARCHt{}. The E1(c) range of the \WC{} statistic
over seeds ($-0.273$ to $0.440$) is from \texttt{e1/e1\_nonstat\_raw.csv}.

\emph{Floor of the calibrated MMD p-values} (\cref{sec:emp-perticker}). In
\texttt{e3/e3\_perticker\_null\_raw.csv}, 31 of the 1000 MMD null p-values equal the smallest
attainable wild-bootstrap value $1/501$, so the smallest calibrated p-value is $32/1001\approx0.032$.
The Benjamini--Hochberg procedure at 5\% over 100 tickers rejects the $k$ smallest p-values for
the largest $k$ with $p_{(k)}\le0.05k/100$; with $p_{(k)}\ge0.032$ this requires $k\ge64$.

\emph{Observed $\pmed$ against the boundary replications} (\cref{sec:emp-verdicts}). For each
replication of the power check we recompute the rule's $\pmed$ from the five single-seed
boundary p-values (\texttt{e3/e3\_realvol\_raw.csv} for D2, \texttt{e3/e3\_power\_raw.csv} for D1)
and count the replications with $\pmed$ at most the observed value of each generator. In D2 (100
replications) the counts are 1 for GBM, 1 for Merton, 4 for Heston and 36 for \GARCHt{}; in D1 they
are 0, 0, 0 and 19. The corresponding Monte Carlo p-values $(1+\text{count})/101$ are 0.020,
0.020, 0.050 and 0.366 in D2. For the bootstrap (13 in D2) the comparison is not meaningful,
because its anchored test is invalid.

\emph{Holm-adjusted p-values} (\cref{sec:emp-verdicts}). Holm's step-down adjustment over the four
parametric generators, computed by the script \texttt{e3\_verdicts\_cal\_main.py} (in
\texttt{paper/scripts}) and recorded in the comment lines of the table it writes: for
the nominal anchored $\pmed$, 0.002 (GBM), 0.016 (Merton), 0.033 (Heston) and 0.432 (\GARCHt{});
for the D2-calibrated p-values, 0.079 for GBM and Merton, 0.099 for Heston and 0.366 for
\GARCHt{}.

\emph{Fixed realised-volatility statistic} (\cref{sec:emp-witness}). \Cref{tab:e3-rv-baseline}
uses the seed-0 evaluation windows stored in \texttt{e3/e3\_witness.csv}, for the absolute test
only; it is written by the script \texttt{e3\_rv\_baseline.py} in \texttt{paper/scripts}. The
script recomputes the critic statistic $T$ and its $\zhac$ from the stored critic values, checks
them against the seed-0 values in \texttt{e3/e3\_pooled\_raw.csv}, and then replaces the critic by
$h(w)=-\log\mathrm{RV}(w)$ on the same windows, with the same Driscoll--Kraay variance.

\begin{table}[tbp]
  \centering
  \caption{What the seed-0 \WC{} critics respond to (E3, $L=20$; 1500 real and 1500 synthetic
    evaluation windows per generator). $T$: the seed-0 statistic. Spearman correlations of the
    critic value with the window log realised volatility and with the largest absolute
    normalised return of the window, over the 3000 windows. Share of $T$ from realised
    volatility: the critic replaced by its mean within 20 equal-count bins of log realised
    volatility. Median realised volatility (root mean square of the normalised returns of a
    window; 1 = \TRAIN{} standard deviation of the ticker) of the real and synthetic evaluation
    windows. Derived from \texttt{results/e3/e3\_witness.csv}.}
  \label{tab:e3-witness-app}
  \small
  \fittable{manual_e3_witness_app}
\end{table}
